\chapter{Active Circuits, Feedback, and Oscillators}
\label{ch:active}
Everything so far has been \emph{passive}: circuits that store and dissipate the
energy handed to them. An active circuit adds a device that can supply gain, and with
gain comes the one operation passive circuits cannot perform---feedback around a
loop. This chapter is therefore where \cref{ch:feedback} finally meets hardware, and
where the resonators of \cref{ch:rlc,ch:rlcpar} stop being things to analyze and
become things a designer places deliberately.

\section{Why Active Circuits Appear}
The Extra exam includes amplifiers, op-amps, feedback, oscillators, mixers, and
regulated power supplies. A full electronics text is not needed here; the goal is
to connect the circuit formulas to system behavior---and, in this chapter, to the
classical feedback ideas that a controls background already contains. The device
details (bias points, transistor models, tube parameters) live in the ARRL
manual; what follows is the systems view that ties them together. The feedback
groundwork of \cref{ch:feedback}---closed-loop gain \(A/(1+A\beta)\), the
large-loop-gain limit \(A_{\text{cl}}\to1/\beta\), pole motion under feedback, and
the gain and phase margins---is the backbone of this chapter. Rather than restate
those results, this chapter puts them to work on real devices, so it is worth having
\cref{ch:feedback} fresh in mind.

\section{Amplifier Gain and Classes}
\label{sec:classes}
Voltage gain is \(A_v=V_o/V_i\); in decibels \(A_{v,\dB}=20\log_{10}|A_v|\), and power
gain is \(G_{\dB}=10\log_{10}(P_o/P_i)\)---the \(20\log\) versus \(10\log\)
distinction of \cref{sec:decibels}.

Amplifier \emph{class} describes how much of the input cycle the active device
conducts, which trades linearity against efficiency:
\begin{description}
\item[Class A] conducts the full \(360^\circ\); most linear, least efficient
(theoretical max \SI{50}{\percent} transformer-coupled, near \SI{25}{\percent}
resistively coupled).
\item[Class B / AB] Class B conducts half the cycle (\(180^\circ\)); Class AB
somewhat more than half. Higher efficiency than Class A, with some crossover
distortion; common in linear single-sideband (SSB) power amplifiers.
\item[Class C] conducts less than half the cycle; efficient but nonlinear, used
where a tuned output tank restores the sinusoid (as in frequency-modulation, FM,
and continuous-wave Morse, CW, transmitters).
\end{description}
Class C is nonlinear on purpose: the output tank is the high-\(Q\) parallel resonator
of \cref{ch:rlcpar} (a pole pair near the \(\jj\omega\) axis) that rings at the
fundamental and rejects the harmonics the clipped current injects. The same nonlinearity that spoils
linearity is what makes mixers and frequency multipliers work.

\section{Op-Amp Ideals}
An ideal op-amp has very high input impedance, very low output impedance, and
very large open-loop gain \(A(s)\).

\begin{physicalbox}
\textbf{How the op-amp is modeled.} An op-amp is \emph{not} a reactive element
like an inductor or capacitor---it stores no energy and contributes no pole of its
own in the ideal picture. It is modeled as a \textbf{dependent source}:
specifically a voltage-controlled voltage source (VCVS) whose output is
\(A(s)\,(v_+-v_-)\), the open-loop gain times the difference of the two input
voltages. The idealizations attach to that source: infinite input impedance means
no current flows into the \(\pm\) terminals; zero output impedance means the
controlled source drives the output stiffly; and \(A(s)\) is enormous (and rolls
off with frequency, \cref{sec:gbw}). Every op-amp result below is bookkeeping on
this one dependent-source model plus the feedback wired around it.
\end{physicalbox}

Negative feedback makes practical closed-loop gains predictable:
\[
\text{inverting:}\quad \frac{V_o}{V_i}=-\frac{R_f}{R_{in}},
\qquad
\text{non-inverting:}\quad \frac{V_o}{V_i}=1+\frac{R_f}{R_g}.
\]
These come from the two ideal rules---no current into the inputs, and the inputs
driven to equal voltage by feedback---but both rules are consequences of a deeper
statement: the loop gain is large. The inverting configuration
(\cref{fig:opamp-inv}) makes this concrete: with the non-inverting input grounded
and the inverting input held at a virtual ground by feedback, the input current
\(V_i/R_{in}\) must flow on through \(R_f\), giving \(V_o=-R_f\,(V_i/R_{in})\).

\subsection{Where \texorpdfstring{\(\beta\)}{beta} Is Hiding}
\label{sec:beta-resistors}
\Cref{ch:feedback} worked entirely with an abstract feedback fraction \(\beta\), and
it is worth connecting that symbol to actual components before going further---because
every result in that chapter then becomes an op-amp result.

Look at the non-inverting amplifier. The output drives the divider \(R_f\)--\(R_g\),
and the voltage fed back to the inverting input is what that divider delivers:
\[
v_-=V_o\,\frac{R_g}{R_f+R_g}
\qquad\Longrightarrow\qquad
\boxed{\;\beta=\frac{R_g}{R_f+R_g}\;}
\]
So \(\beta\) is nothing more exotic than \cref{ch:resistive}'s voltage divider,
running backwards from output to input. The large-loop-gain limit
\(A_{\mathrm{cl}}\to1/\beta\) of \cref{ch:feedback} then reads
\[
\frac{1}{\beta}=\frac{R_f+R_g}{R_g}=1+\frac{R_f}{R_g},
\]
which is exactly the non-inverting gain formula above. The textbook result and the
control-theory result are the same statement.

Keeping \(A\) finite shows what the ideal rules cost. Substituting into
\(A_{\mathrm{cl}}=A/(1+A\beta)\) and factoring out \(1/\beta\),
\[
A_{\mathrm{cl}}=\frac{1}{\beta}\cdot\frac{A\beta}{1+A\beta}
=\underbrace{\Bigl(1+\frac{R_f}{R_g}\Bigr)}_{\text{ideal}}
\cdot\underbrace{\frac{1}{1+1/L}}_{\text{error}},
\qquad L=A\beta ,
\]
so the fractional gain error is about \(1/L\)---\cref{ch:feedback}'s sensitivity
result, in the form a designer uses. With \(A=10^5\) and a gain of \(100\)
(\(\beta=10^{-2}\), so \(L=1000\)) the error is \(0.1\%\), which is why the ideal
formula is safe.

The same substitution quantifies the virtual ground. The op-amp's differential input
voltage is not zero but
\[
v_+-v_-=\frac{V_o}{A},
\]
so for \(V_o=\SI{10}{\volt}\) and \(A=10^5\) the two inputs differ by
\SI{100}{\micro\volt}. ``The inputs are at the same voltage'' is not an axiom---it is
a consequence of \(A\) being enormous, and it degrades exactly as \(A\) falls with
frequency.

\begin{figure}[htbp]
\centering
\begin{circuitikz}[scale=1]
\node[op amp,yscale=-1] (oa) at (0,0) {};
\draw (oa.-) to[short] ++(-0.7,0) coordinate(n);
\draw (n) to[R,l_=\(R_{in}\)] ++(-2.2,0) node[left]{\(V_i\)};
\draw (oa.+) to[short] ++(-0.7,0) node[ground]{};
\draw (oa.out) to[short,-o] ++(0.9,0) coordinate(vo) node[right]{\(V_o\)};
\draw (n) -- ++(0,1.6) coordinate(tl) to[R,l=\(R_f\)] (tl-|vo) -- (vo);
\end{circuitikz}
\caption{Inverting op-amp amplifier. Feedback holds the inverting input at a
virtual ground, so \(V_o/V_i=-R_f/R_{in}\)---set by the resistor ratio, not the
op-amp.}
\label{fig:opamp-inv}
\end{figure}

\section{Feedback as a Loop}
Model the amplifier as a forward path \(A(s)\) with a fraction \(\beta\) of the
output fed back and subtracted at the input (\cref{fig:feedback-loop}). This is
\cref{ch:feedback}'s block diagram with the amplifier's names attached, so the
results below are that chapter's results read as hardware.

\begin{figure}[htbp]
\centering
\begin{tikzpicture}[>=Latex,node distance=1.1cm,scale=1]
\node (in) at (0,0) {\(V_i\)};
\node[draw,circle,inner sep=1.4pt] (sum) at (1.3,0) {\(\Sigma\)};
\node[draw,minimum width=1.3cm,minimum height=0.8cm] (A) at (3.6,0) {\(A(s)\)};
\node (out) at (6.3,0) {\(V_o\)};
\node[draw,minimum width=1.1cm,minimum height=0.7cm] (B) at (3.6,-1.5) {\(\beta\)};
\draw[->] (in) -- (sum);
\draw[->] (sum) -- (A);
\draw[->] (A) -- (out);
\draw[->] (5.3,0) |- (B);
\draw[->] (B) -| (sum);
\node at (0.95,0.3) {\small \(+\)};
\node at (1.62,-0.28) {\small \(-\)};
\end{tikzpicture}
\caption{Standard negative-feedback loop: the reference \(V_i\) enters the summing
junction with \(+\) and the feedback \(\beta V_o\) is subtracted (\(-\)). The loop
gain is \(L(s)=A(s)\beta\).}
\label{fig:feedback-loop}
\end{figure}

The closed-loop gain is
\[
A_{\text{cl}}(s)=\frac{A(s)}{1+A(s)\beta}
=\frac{A(s)}{1+L(s)},
\qquad
L(s)=A(s)\beta .
\]
When the loop gain \(|L|\gg1\), \(A_{\text{cl}}\to 1/\beta\): the gain is set by
the passive feedback network, not by the messy device parameters. That single
fact is why the inverting and non-inverting formulas above depend only on
resistor ratios.

\begin{controlsbox}
\textbf{What this buys.} By the sensitivity result of \cref{ch:feedback},
\(\dd A_{\text{cl}}/A_{\text{cl}}=(\dd A/A)/(1+L)\): the loop divides the device's
imperfections by \(1+L\). For an op-amp that is the whole design philosophy---spend
the enormous surplus gain \(A\) to buy immunity from device variation, temperature,
and distortion, and let two resistors set the answer.
\end{controlsbox}

\section{Gain--Bandwidth Product}
\label{sec:gbw}
An internally compensated op-amp is deliberately built so that one pole dominates,
\[
A(s)=\frac{A_0}{1+s/\omega_a},
\]
with an enormous DC gain \(A_0\) (often \(10^5\)) and a very low open-loop corner
\(\omega_a\) (often a few hertz). That is precisely the first-order plant analyzed
in \cref{ch:feedback}, so its result transfers unchanged: closing the loop with
constant \(\beta\) divides the DC gain by \((1+A_0\beta)\) and multiplies the corner
by the same factor, leaving the product invariant,
\[
(\text{gain})\times(\text{bandwidth})=A_0\,\omega_a=\text{constant}.
\]
For amplifiers this invariant has a name and a datasheet value: the
\textbf{gain--bandwidth product}, quoted as the \emph{unity-gain frequency}
\(f_T=A_0f_a\)---the frequency at which the op-amp can no longer provide any gain at
all. Every closed-loop configuration you build with that part lands on the same
diagonal (\cref{fig:opamp-gbw}), so choosing a closed-loop gain \emph{is} choosing a
bandwidth: doubling one halves the other.

\begin{figure}[htbp]
\centering
\includegraphics[width=0.86\textwidth]{opamp_gbw_bode}
\caption{Gain--bandwidth product of a single-pole op-amp with
\(A_0=10^5\) (\SI{100}{\decibel}) and \(f_a=\SI{10}{\hertz}\), so
\(f_T=\SI{1}{\mega\hertz}\). Each closed-loop gain (set by \(1/\beta\)) is flat until
it meets the open-loop skirt, then rolls off along it, so every closed-loop corner
lands on the same unity-gain frequency \(f_T=A_0f_a\). Trading gain for bandwidth
slides along one fixed diagonal.}
\label{fig:opamp-gbw}
\end{figure}

\begin{workedbox}[: reading a real part's limits]
The op-amp plotted in \cref{fig:opamp-gbw} is a common one: \(A_0=10^5\)
(\SI{100}{\decibel}) with an open-loop corner at \(f_a=\SI{10}{\hertz}\), so
\[
f_T=A_0f_a=(10^5)(\SI{10}{\hertz})=\SI{1}{\mega\hertz}.
\]
Build a non-inverting amplifier with a gain of \(100\). From
\cref{sec:beta-resistors}, that needs \(1+R_f/R_g=100\)---say
\(R_f=\SI{99}{\kilo\ohm}\) and \(R_g=\SI{1}{\kilo\ohm}\), giving
\(\beta=1/100\). Then:
\[
L(0)=A_0\beta=1000,
\qquad
A_{\mathrm{cl}}(0)=\frac{10^5}{1+1000}=99.9,
\qquad
f_c=f_a(1+A_0\beta)=\SI{10.0}{\kilo\hertz}.
\]
Three things to read off. The gain is \(99.9\) rather than \(100\)---the \(0.1\%\)
error predicted by \(1/L\). The bandwidth is \SI{10}{\kilo\hertz}, not
\SI{1}{\mega\hertz}: gain \(\times\) bandwidth returns \(f_T\), as it must. And the
closed-loop pole has moved from \SI{10}{\hertz} out to \SI{10}{\kilo\hertz}, a factor
of \(1001\) to the left---\cref{ch:feedback}'s first-order result with real numbers.
\end{workedbox}

\begin{exambox}
This is the practical shape of gain--bandwidth questions. An op-amp with
\(f_T=\SI{1}{\mega\hertz}\) configured for a gain of \(100\) gives about
\SI{10}{\kilo\hertz} of bandwidth---fine for audio, useless at RF. A stage needing
both high gain and wide bandwidth needs either a faster part or two cascaded stages of
modest gain (whose gains multiply while each keeps its own wider bandwidth).
\end{exambox}

As \cref{ch:feedback} showed, this single-pole model can never be destabilized by
more feedback---the pole only slides left. Instability requires the extra phase lag
of additional poles, which every real amplifier has.

\section{Stability, Phase Margin, and Gain Margin}
Real amplifiers have more than one pole. With two or more poles the loop gain
\(L(\jj\omega)\) accumulates phase lag, and if the phase reaches \(-180^\circ\)
while \(|L|\ge 1\), the subtracting junction effectively adds---the feedback has
become positive at that frequency and the closed-loop poles move toward, or
into, the right half-plane. \Cref{sec:threepoles} counts how many poles that takes:
one or two can never satisfy both conditions at once, three can, and the same
phase accounting fixes the critical gain at \(K=8\) for three identical lags.

The gain and phase margins of \cref{ch:feedback} are exactly how amplifier
stability is specified and measured: positive margins keep the closed-loop poles in
the left half-plane, small margins put them near the \(\jj\omega\) axis (ringing,
overshoot, a peaked response), and zero margin puts them on it. An op-amp datasheet's
phase-margin figure is that number for the part's own internal loop.

The same story reads on the s-plane as a root locus (\cref{fig:opamp-rootlocus}):
as loop gain rises, the closed-loop poles of a multi-pole amplifier migrate, and with
three poles two of them curve toward and cross the \(\jj\omega\) axis---the moment
small margins become ringing and then sustained oscillation.

\begin{controlsbox}
\textbf{One amplifier, three figures.} \Cref{fig:opamp-rootlocus} is deliberately the
\emph{same} three-pole loop \(L(s)=K/(1+s/\omega_0)^3\) whose Bode margins appear in
\cref{fig:fb-margins} and whose oscillation condition appears in
\cref{fig:oscillator-loopgain}. Follow one amplifier across the three:
\begin{itemize}
\item At \(K=3\), \cref{fig:fb-margins} shows healthy gain and phase margins, and the
locus has its poles well inside the left half-plane.
\item Raising \(K\) walks the poles right along the locus, shrinking both margins.
\item At \(K=8\) the poles reach the \(\jj\omega\) axis at
\(\omega=\sqrt3\,\omega_0\); \cref{fig:oscillator-loopgain} shows the margins hitting
exactly zero, which is \(L=-1\).
\item Beyond \(K=8\) the poles are in the right half-plane and the amplifier
oscillates whether you wanted it to or not.
\end{itemize}
Margins and root locus are not two topics---they are two readings of this one
parameter sweep, which is the claim \cref{ch:feedback} opened with.
\end{controlsbox}

\begin{figure}[htbp]
\centering
\includegraphics[width=0.6\textwidth]{opamp_rootlocus}
\caption{Root locus of the three-pole amplifier \(L(s)=K/(1+s/\omega_0)^3\) as loop
gain \(K\) increases---the same loop whose Bode margins appear in
\cref{fig:fb-margins} and whose oscillation condition appears in
\cref{fig:oscillator-loopgain}. Two closed-loop poles leave the real axis and reach
the \(\jj\omega\) axis at \(K=8\), \(\omega=\sqrt3\,\omega_0\); beyond that the
amplifier oscillates. A single-pole amplifier (one branch on the real axis) can never
do this---instability needs the extra phase lag of additional poles.}
\label{fig:opamp-rootlocus}
\end{figure}

Compensation (deliberately lowering the first pole, as in an internally
compensated op-amp) restores phase margin by making the loop gain fall through
\SI{0}{\decibel} before the second pole adds its lag. It buys stability with
bandwidth---the same trade, again.

\subsection{Neutralization: Cancelling a Loop You Did Not Ask For}
\label{sec:neutralization}
Compensation assumes you know where the loop is. RF amplifiers have a harder
problem: a feedback path exists whether or not the designer put one there.

Every active device has capacitance between its output and its input---collector to
base, drain to gate, plate to grid. That capacitance is an impedance, so it feeds a
fraction of the output back, and at RF the fraction is not small: the feedback
admittance \(\jj\omega C\) grows with frequency, so a device that is unconditionally
stable at audio can have substantial loop gain at \SI{30}{\mega\hertz}. Worse, the
feedback arrives through a \emph{capacitor}, which contributes its own
\(+90^\circ\)---so the phase condition of \cref{sec:threepoles} can be met by the
device's own strays and the stage oscillates at a frequency nobody designed.

\textbf{Neutralization} is the cure, and it is the only place in this book where
feedback is deliberately cancelled rather than shaped. A second path is added from
output to input, matched in magnitude to the stray but opposite in phase---typically
by taking the feedback from the other end of a center-tapped tank, which supplies a
signal \(180^\circ\) out of phase, through an adjustable capacitor set to match the
offending one. The two paths then sum to zero at the input:
\[
\underbrace{\jj\omega C_{\text{stray}}}_{\text{unwanted path}}
+\underbrace{\bigl(-\jj\omega C_{\text{neut}}\bigr)}_{\text{added path}}=0
\qquad\text{when}\qquad
C_{\text{neut}}=C_{\text{stray}} .
\]
With the loop gain driven to zero there is no \(L\) to satisfy \(L=-1\), and the
stage cannot oscillate at any drive level. Screen-grid tubes attack the same problem
structurally rather than by cancellation---the screen is an earthed electrode
interposed between grid and plate, which reduces \(C_{\text{stray}}\) itself---and
modern practice often avoids the issue by choosing devices whose feedback capacitance
is small enough to leave the loop gain below unity throughout.

\begin{controlsbox}
\textbf{Three ways to keep a loop stable, and they are not interchangeable.}
Compensation reshapes \(L(\jj\omega)\) so its magnitude falls below unity before its
phase reaches \(-180^\circ\); it keeps the loop and makes it well behaved.
Neutralization \emph{removes} the loop by summing it against an equal and opposite
copy. Choosing a better device shrinks the loop's gain at source. All three are
statements about the same \(L\), which is why one framework covers an op-amp, a
regulator, and a neutralized RF power amplifier.
\end{controlsbox}

\section{Oscillators: Poles Placed on the Axis}
An oscillator is a feedback loop designed to sit exactly on the stability
boundary. The closed-loop poles are the roots of \(1+L(s)=0\), so they reach the
\(\jj\omega\) axis when
\[
L(\jj\omega_{\text{osc}})=-1,
\qquad\text{that is}\qquad
|L(\jj\omega_{\text{osc}})|=1,\quad
\angle L(\jj\omega_{\text{osc}})=-180^\circ,
\]
the Nyquist \(-1\) point of \cref{sec:nyquist}---the oscillator is designed so its
loop-gain curve passes exactly \emph{through} that point instead of clearing it.
This is the \emph{Barkhausen criterion}. Its familiar
``unit magnitude, zero net phase'' wording refers to the round-trip transmission
\(T=-L\), which already includes the \(-1\) of the subtracting junction; setting
\(T=1\) is the same statement as \(L=-1\). In s-plane language the closed-loop
poles are placed on the \(\jj\omega\) axis at \(\pm\jj\omega_{\text{osc}}\)---the
marginally stable case from the s-plane chapter. At start-up the design makes \(|L|>1\) so the pole pair sits slightly in
the right half-plane and the amplitude grows; a nonlinearity (device saturation,
a limiter, a lamp, or automatic gain control, AGC) then reduces the effective gain until \(|L|=1\) and the
poles settle onto the axis at a stable amplitude.

\Cref{fig:oscillator-loopgain} shows the condition on a loop-gain Bode plot: the
design places the frequency where the phase reaches \(-180^\circ\) exactly at the
\(\SI{0}{\decibel}\) (\(|L|=1\)) crossing, so \(L=-1\) is met at one frequency.

\begin{figure}[htbp]
\centering
\includegraphics[width=0.82\textwidth]{oscillator_loopgain}
\caption{Barkhausen condition for a three-pole loop
\(L(s)=K/(1+s/\omega_0)^3\) with \(K=8\). The magnitude passes through
\SI{0}{\decibel} at the same frequency \(\omega_{\mathrm{osc}}=\sqrt3\,\omega_0\)
where the phase reaches \(-180^\circ\): loop gain unity with \(180^\circ\) of lag
is \(L=-1\), so the closed-loop poles sit on the \(\jj\omega\) axis.}
\label{fig:oscillator-loopgain}
\end{figure}

The frequency-setting network sets where those poles land:
\begin{description}
\item[LC oscillators] (Colpitts, Hartley) use an LC tank; \(\omega_{\text{osc}}
\approx 1/\sqrt{LC}\). Higher tank \(Q\) means poles closer to the ideal axis
location and lower phase noise.
\item[Crystal oscillators] replace the tank with a quartz resonator of very high
\(Q\), giving excellent frequency stability.
\item[RC phase-shift / Wien-bridge] use RC networks to reach the required loop
phase; common at audio frequencies.
\end{description}

\begin{physicalbox}
Phase noise is the frequency-domain fingerprint of how firmly the poles are
pinned to the axis. A high-\(Q\) resonator holds the pole pair tightly at
\(\pm\jj\omega_{\text{osc}}\), so the spectrum is a sharp line; a low-\(Q\) loop
lets the poles wander, smearing the line. This is why crystal and high-\(Q\) LC
oscillators are quieter.

\Cref{sec:phasenoise} puts a number on it: the skirt inherits the tank's
single-pole roll-off, so phase noise falls at \(-\SI{20}{\decibel}\) per decade of
offset and scales as \(1/Q_L^2\)---\emph{doubling the loaded \(Q\) buys
\SI{6}{\decibel}}. It also shows what that costs a receiver, since a noisy local
oscillator raises the apparent noise floor near strong signals.
\end{physicalbox}

\section{Active Filters and Their Limits}
Op-amps can implement active RC filters that synthesize the second-order pole
pairs of \cref{ch:filters} without inductors, by using feedback to create complex
poles from resistors and capacitors. That claim deserves scrutiny: two RC stages in
cascade give two \emph{real} poles (\cref{ch:rc}), and no arrangement of resistors and
capacitors alone can produce a complex pair. Feedback is what does it.

The Sallen--Key low-pass (\cref{fig:sallenkey}) is the classic example. Writing KCL at
the two RC nodes with the buffer enforcing \(v_-=V_o\)---so that \(C_2\) feeds the
output back to the \emph{first} node rather than to ground---and eliminating the middle
node gives
\[
G(s)=\frac{\omega_0^2}{s^2+\dfrac{\omega_0}{Q}s+\omega_0^2},
\qquad
\omega_0=\frac{1}{\sqrt{R_1R_2C_1C_2}},
\qquad
\frac{1}{Q}=\sqrt{\frac{C_2}{C_1}}
\left(\sqrt{\frac{R_1}{R_2}}+\sqrt{\frac{R_2}{R_1}}\right).
\]
This is precisely the standard second-order form of \cref{ch:linsys}---the same one the
LC circuit of \cref{ch:rlc} produced---so everything about \(\omega_0\), \(\zeta\), and
resonant peaking transfers unchanged. That is also what makes one stage enough to
build any order: by \cref{sec:factoring} a filter of order \(n\) is a product of
sections of exactly this form, and \cref{sec:poleplacement} tabulates the
\((\omega_0,Q)\) each section needs. An order-5 Butterworth, for instance, is two of
the stages below---one at \(Q=0.618\), one at \(Q=1.618\)---plus a single RC section
for the leftover real pole. What is new is the \emph{knob}: with equal
resistors the bracket becomes \(2\), leaving
\[
Q=\frac{1}{2}\sqrt{\frac{C_1}{C_2}} ,
\]
so \textbf{the capacitor ratio sets the damping}. That is the whole trick. In an LC
filter \(Q\) is set by resistance (\cref{sec:lc-lowpass}); here it is set by a ratio of
two capacitors, and the feedback path through \(C_2\) is what makes the pole pair
complex. No inductor appears anywhere.

\begin{workedbox}[: a Butterworth audio section]
Design a second-order Butterworth low-pass at \(f_0=\SI{3}{\kilo\hertz}\).
Butterworth means \(Q=1/\sqrt2\approx0.707\) (\cref{ch:filters}), so with equal
resistors,
\[
Q=\frac{1}{2}\sqrt{\frac{C_1}{C_2}}=\frac{1}{\sqrt2}
\qquad\Longrightarrow\qquad
C_1=2\,C_2 .
\]
Choosing \(R_1=R_2=\SI{10}{\kilo\ohm}\) and imposing
\(\omega_0=1/(R\sqrt{C_1C_2})=2\pi(\SI{3}{\kilo\hertz})\) gives
\[
C_2\approx\SI{3.75}{\nano\farad},
\qquad
C_1\approx\SI{7.5}{\nano\farad}.
\]
Two resistors, two capacitors, one op-amp---and a maximally flat second-order response
that would otherwise need an inductor. Check the limits, though: at
\(f_0=\SI{3}{\kilo\hertz}\) an op-amp with \(f_T=\SI{1}{\mega\hertz}\) still has a loop
gain of roughly \(300\), so the ideal analysis holds comfortably. Push the same
topology to \SI{300}{\kilo\hertz} and it would not.
\end{workedbox}

\begin{figure}[htbp]
\centering
\begin{circuitikz}[scale=1]
\node[op amp] (oa) at (7,0) {};
\draw (0,0) node[left]{\(V_{\mathrm{in}}\)} to[R,l=\(R_1\)] (2,0) coordinate(a)
      to[R,l=\(R_2\)] (4,0) coordinate(b);
\draw (b) -| (oa.+);
\draw (b) to[C,l_=\(C_1\)] ($(b)+(0,-1.8)$) node[ground]{};
\draw (oa.out) to[short,-o] (9,0) coordinate(vo) node[right]{\(V_o\)};
% unity-gain feedback from inverting input up and over to output
\draw (oa.-) -- ++(0,1.2) -| (vo);
% C2 from node a up and over to the output
\draw (a) -- ++(0,2.5) coordinate(ct) to[C,l=\(C_2\)] (ct-|vo) -- (vo);
\end{circuitikz}
\caption{Sallen--Key second-order low-pass. Two RC stages around a unity-gain
op-amp buffer create a complex pole pair without an inductor;
\(\omega_0=1/\sqrt{R_1R_2C_1C_2}\).}
\label{fig:sallenkey}
\end{figure}

The catch is that the same finite gain--bandwidth product limits them: near the
amplifier's own gain-crossover the ideal-op-amp assumption fails, so an active
filter that behaves at audio may be useless at RF, where LC or crystal filters
take over.

\section{Mixers: Where Linearity Ends}
\label{sec:mixers}
Every analytical tool in this book---transfer functions, poles, Bode plots, root
loci---rests on one assumption stated back in \cref{ch:linsys}: the system is
\emph{linear}. A mixer is the one circuit in radio built to violate it deliberately,
and saying so plainly is worth more than papering over it.

A linear system obeys superposition: it can change a sinusoid's amplitude and phase,
but it can never change its \emph{frequency}. That is precisely why \(G(\jj\omega)\)
is meaningful---one input frequency in, the same frequency out. A mixer must produce
frequencies that were not present at its input, so no transfer function can describe
it. What it does instead is \emph{multiply}:
\[
v_{\text{out}}\propto
\underbrace{\cos(\omega_{\text{RF}}t)}_{\text{signal}}\times
\underbrace{\cos(\omega_{\text{LO}}t)}_{\text{local oscillator}}
=\tfrac12\cos\bigl[(\omega_{\text{LO}}-\omega_{\text{RF}})t\bigr]
+\tfrac12\cos\bigl[(\omega_{\text{LO}}+\omega_{\text{RF}})t\bigr],
\]
using the product-to-sum identity. Two inputs, and the output contains their
\textbf{difference} and \textbf{sum} frequencies---neither of which was there before.
This is the whole basis of the superheterodyne receiver: multiply the incoming signal
by a local oscillator, keep the difference, and you have translated the signal to a
fixed intermediate frequency where the sharp filters of \cref{ch:filters} and the
high-gain amplifiers of this chapter can do their work at one frequency forever.

The same multiplication has one more use, and it is the reason \cref{ch:dsp} needs
almost no new theory: \emph{sampling} a signal is multiplying it by a train of
impulses, which is multiplication by every multiple of the sampling rate at once
(\cref{sec:samplingismult}). A mixer makes two frequency-shifted copies; a sampler
makes infinitely many.

Two consequences follow immediately from the algebra.

\emph{Image response.} The difference term depends only on
\(|\omega_{\text{LO}}-\omega_{\text{RF}}|\), so \emph{two} input frequencies---one
above the local oscillator and one below, equally spaced---produce the same
intermediate frequency. The unwanted one is the \textbf{image}, and no amount of
filtering \emph{after} the mixer can separate it from the wanted signal, because by
then they are identical. It must be rejected before mixing, which is what the front-end
filter is for.

The image's location follows from the same absolute value, and is worth having as a
number. If the wanted signal sits at \(f_{\text{LO}}-f_{\text{IF}}\), then the
frequency \(f_{\text{LO}}+f_{\text{IF}}\) mixes to the same place, so the two are
separated by
\[
f_{\text{image}}-f_{\text{signal}}=2f_{\text{IF}} .
\]
\emph{The image is always twice the intermediate frequency away.} That single fact
drives receiver architecture: a high first IF pushes the image far out where a
modest front-end filter can reject it, which is exactly why up-converting receivers
use first IFs above the highest received frequency.

\subsection{The Same Multiplier, Four Jobs}
A mixer's algebra does not care which port you call which, so the same circuit
appears under four names depending on what is connected and what is kept:
\begin{description}
\item[Mixer / converter] RF in, local oscillator in, keep the difference: translate
a signal to the IF. The receiving case above.
\item[Balanced modulator] audio in, carrier in, and---by balancing the circuit so
the two inputs cancel at the output---keep \emph{only} the sum and difference. The
carrier itself is suppressed, leaving double sideband (DSB). Filter off one of the
two sidebands and the result is single sideband; this is how an SSB transmitter
generates its signal.
\item[Product detector] IF in, reinserted carrier in, keep the difference: the
audio. It is the balanced modulator run backwards, which is why an SSB receiver and
an SSB transmitter contain the same part twice.
\item[Frequency multiplier] one input, driven hard into nonlinearity, and keep a
\emph{harmonic} rather than a mixing product---the \(m f_1\) terms with a tuned
circuit selecting one.
\end{description}
Recognizing that these are one operation with different ports and different
post-filters is worth more than memorizing four block diagrams.

\emph{Intermodulation.} Real mixers and overdriven amplifiers do not multiply cleanly;
their nonlinearity contains higher-order terms, which generate combinations
\(m f_1\pm n f_2\) of any two signals present. The third-order products
\(2f_1-f_2\) and \(2f_2-f_1\) are the troublesome ones, because for two nearby
signals they land \emph{nearby too}---inside the passband, where no filter can remove
them. That is why receiver front ends are specified for intermodulation
performance, and why an amplifier driven into nonlinearity produces spurious signals
rather than merely a distorted one.

\Cref{sec:powerseries} makes that nonlinearity explicit---expanding the device as
\(y=a_1x+a_2x^2+a_3x^3\) and driving it with two tones---and shows that every
strong-signal specification a receiver carries is a consequence of those three
coefficients.

Concretely, with \(f_1=\SI{14.20}{\mega\hertz}\) and \(f_2=\SI{14.21}{\mega\hertz}\),
the third-order products land at
\[
2f_1-f_2=\SI{14.19}{\mega\hertz},
\qquad
2f_2-f_1=\SI{14.22}{\mega\hertz},
\]
\SI{10}{\kilo\hertz} either side of the pair---the same spacing as the originals,
and hopelessly inside any filter wide enough to pass them. Even-order products, by
contrast, land near DC or near \(2f\), where filtering removes them easily. That
asymmetry is why odd-order intermodulation is the specification that matters.

\begin{physicalbox}
\textbf{Why the two-tone test is the standard measurement.} The algebra above also
prescribes how to measure linearity. Feed the amplifier two equal tones that are
\emph{not} harmonically related---so that nothing they produce can be confused with
simple harmonic distortion---and look at the output spectrum. A perfectly linear
stage returns exactly two lines. Any third-order curvature announces itself as a
symmetric pair of new lines at \(2f_1-f_2\) and \(2f_2-f_1\), whose level relative
to the wanted tones \emph{is} the intermodulation figure. Two tones rather than one,
because a single tone would produce only harmonics---far away, easily filtered, and
therefore uninformative about the in-band products that actually matter.
\end{physicalbox}

\begin{physicalbox}
\textbf{The same nonlinearity, wanted and unwanted.} Class C amplification
(\cref{sec:classes}) and frequency multiplication rely on exactly the mechanism that
makes intermodulation a problem: a nonlinear device generating harmonics and mixing
products. The difference is only whether a tuned circuit afterwards keeps the product
you wanted. A high-\(Q\) tank (\cref{ch:rlcpar}) keeps one and discards the rest,
which is why Class C works at all---the nonlinearity makes a mess, and a parallel
resonator cleans it up.
\end{physicalbox}

\begin{exambox}
When a question involves two frequencies producing a third, stop looking for a
transfer function---the circuit is nonlinear and the answer is arithmetic on
frequencies: sums, differences, and \(mf_1\pm nf_2\). When a question involves one
frequency in and the same frequency out with altered amplitude or phase, the circuit
is linear and every tool in Part~II applies.
\end{exambox}

\section{Power Supplies}
A power supply looks like a different subject and is not: it is a rectifier, an RC
discharge, a low-pass filter, and a feedback loop, all of which are already in hand.
Taking it in that order costs nothing new.

\subsection{Rectification and Ripple}
A \textbf{rectifier} uses diodes to make current flow one way only. Two
arrangements matter, and the difference between them is a conduction angle:
\begin{description}
\item[Half-wave] one diode, which conducts on one polarity and blocks the other, so
it conducts through \(180^\circ\) of each input cycle. The output is a train of
pulses at the \emph{line} frequency \(f_{\text{line}}\).
\item[Full-wave] two diodes with a center-tapped transformer (or four in a bridge),
arranged so that whichever half-cycle the input is in, some path conducts. Current
therefore flows through the full \(360^\circ\), and---this is the useful
part---\emph{both} half-cycles are folded to the same polarity, so the output pulses
arrive at \(2f_{\text{line}}\).
\end{description}
That doubling is the reason full-wave rectification is worth the extra parts. The
ripple to be removed sits at twice the frequency, so the filter that removes it
needs a corner an octave higher---or, at a fixed corner, attenuates the ripple
\SI{6}{\decibel} harder (\cref{sec:decibels}). From \SI{60}{\hertz} mains, a
half-wave supply must filter \SI{60}{\hertz} and a full-wave supply
\SI{120}{\hertz}.

\subsection{The Filter Is a Low-Pass, and the Bleeder Is an RC Discharge}
What follows the rectifier is a low-pass filter built from exactly the elements of
\cref{ch:filters}: a shunt capacitor, sometimes with a series inductor (a
\emph{choke}) ahead of it to make it second order. Its job is stated in the language
of \cref{sec:lc-lowpass}---pass DC, reject everything at \(2f_{\text{line}}\) and
above---so the design rule is simply that the corner frequency must sit well below
the ripple frequency. A large capacitor is a low corner.

The \textbf{bleeder resistor} across the filter capacitor is
\cref{ch:rc} unchanged. With the supply off, capacitor and bleeder are an isolated
RC pair, so the stored charge decays as \(\ee^{-t/RC}\) and is gone after about
\(5\tau\). This is a safety component: a large filter capacitor at several hundred
volts holds \(E_C=\tfrac12CV^2\) of energy that will otherwise sit there
indefinitely.

\begin{workedbox}[: how long until it is safe?]
A \SI{100}{\micro\farad} filter capacitor charged to \SI{400}{\volt} stores
\[
E_C=\tfrac12CV^2=\tfrac12(100\times10^{-6})(400)^2=\SI{8}{\joule},
\]
which is emphatically dangerous. Fit a \SI{100}{\kilo\ohm} bleeder and
\[
\tau=RC=(10^5)(100\times10^{-6})=\SI{10}{\second},
\]
so the capacitor is effectively discharged after \(5\tau\approx\SI{50}{\second}\).
The bleeder wastes \(V^2/R=\SI{1.6}{\watt}\) continuously in normal operation---the
standard trade, safety against standing loss. Note the two competing pulls: a
smaller resistor discharges faster but wastes more, and the pole \(-1/RC\) is the
only thing being moved.
\end{workedbox}

When several capacitors are stacked in series to withstand a high voltage,
equal-value resistors across each one do double duty: they bleed the stack down, and
because they are a \emph{divider} (\cref{ch:resistive}) they force the total voltage
to share equally instead of according to the capacitors' leakage.

\subsection{Regulation Is a Feedback Loop}
A linear regulator compares its output against a reference and drives a pass element
to null the difference---so it is literally \cref{ch:feedback}'s block diagram with
a transistor as the forward path. Everything from that chapter therefore applies:
the large-loop-gain limit is why the output tracks the reference rather than the
transistor's parameters, the sensitivity result \(1/(1+L)\) is what ``line and load
regulation'' measure, and the margins are why a regulator with too much loop gain
or a marginal output capacitor \emph{rings}. Switching converters replace the
dissipative pass element with a duty-cycle-modulated switch and an LC filter; the
high switching frequency is chosen deliberately, because a higher ripple frequency
means the same attenuation from far smaller reactive components.

\begin{exambox}
Power-supply questions usually reduce to one of three facts. Full-wave rectification
puts ripple at \emph{twice} the line frequency and half-wave at the line frequency.
The filter is a low-pass, so bigger capacitance means a lower corner and less
ripple. And the bleeder is an RC discharge whose \(\tau=RC\) sets how long the supply
stays lethal after switch-off.
\end{exambox}

\begin{exambox}
For active-circuit questions, keep three layers separate: small-signal gain,
power handling, and stability. Feedback questions almost always reduce to one
idea---\emph{loop gain}: large loop gain buys predictable gain, wide bandwidth,
and low distortion; too much phase lag at unity loop gain costs stability; and an
oscillator is the deliberate limiting case with the poles parked on the
\(\jj\omega\) axis.
\end{exambox}
